\chapter{重積分}
    \section{平均値の定理}
        \begin{shadebox}
            $\Omega\in\realNumbers^n$を体積確定の有界閉領域，$f:\Omega\to\realNumbers$を連続関数とすると，$^\exists\bm{x}'\st\integrate{\Omega}{}{f(\bm{x})}{}{\bm{x}} = f(\bm{x}')|\Omega|$
        \end{shadebox}
        \begin{proof}
            \quad\par
            $\Omega$内での$f$の最小値と最大値を各々$m(\Omega),M(\Omega)$とすると
            \[m(\Omega)|\Omega| \leq \integrate{\Omega}{}{f(\bm{x})}{}{\bm{x}} \leq M(\Omega)|\Omega|\quad\therefore\;m(\Omega) \leq \frac{1}{|\Omega|}\integrate{\Omega}{}{f(\bm{x})}{}{\bm{x}} \leq M(\Omega)\]
            であるから中辺の値を$I(\Omega)$とすれば中間値の定理より$^\exists\bm{x}'\st I(\Omega) = f(\bm{x}')$であり，定理が従う。
        \end{proof}
    \section{\texorpdfstring{$C^1$}{C\string^1}級の変換による面積0の集合の像の面積もまた0}
        \begin{shadebox}
            閉超直方体$R\subset\realNumbers^n$を含む開集合$D$上で定義された$C^1$級の\textcolor{red}{変換}$F:D\to\realNumbers^n$を考える。
            集合$C\subset\realNumbers^n$の面積が0であるとき，$F(C)\subset\realNumbers^n$の面積もまた0である。
        \end{shadebox}
        \begin{proof}
            \quad\par
            $F$は$R$で$C^1$級なので特にLipschitz連続である。Lipschitz定数を$K$とする。
            $R = [a_1,b_1]\times\dotsb\times[a_n,b_n]$の$m$等分割$\Delta$を考える。
            第$d$軸方向の分割を$\Delta_d : x^{(d)}_i = a_d + \frac{i}{m}(b_d-a_d)\;(i=0,\dotsc,m,d=1,\dotsc,n)$とし，各小超直方体を$R_{i_1,\dotsc,i_n} \coloneq [x^{(1)}_{i_1-1}, x^{(1)}_{i_1}]\times\dotsb\times[x^{(n)}_{i_n-1}, x^{(n)}_{i_n}]$で表す。
            $R$の対角線の長さを$L = \sqrt{\sum_{d=1}^n (b_d-a_d)^2}$としておく。
            $C$は面積確定で面積が0であるから，任意の$\varepsilon>0$に対してある$M(\varepsilon)\in\naturalNumbers$が存在し，$m\geq M(\varepsilon)$なる任意の分割に対して次式を満たす。
            \[ \frac{|R|}{(2KL)^n}\varepsilon > \sum_{R_{i_1,\dotsc,i_n}\cap C\neq \emptyset} |R_{i_1,\dotsc,i_m}| = \sum_{R_{i_1,\dotsc,i_n}\cap C\neq \emptyset}\frac{|R|}{m^n} \]
            各$R_{i_1,\dotsc,i_n}$の対角線の長さは$L/m$であり，$F$がLipschitz連続だから各$R_{i_1,\dotsc,i_n}$について，その中の任意の1点の$F$による像(点)の半径$KL/m$以内に他の全ての点の像が含まれる。
            つまり各$F(R_{i_1,\dotsc,i_n})$は1辺の長さが$2KL/m$の$n$次元超立方体$R'_{i_1,\dotsc,i_n}$で覆うことができる。
            よって$F(C)$は$\setComprehension*{R'_{i_1,\dotsc,i_n}}{R_{i_1,\dotsc,i_n}\cap C\neq\emptyset}$で覆うことができて，その面積は
            \[ \sum_{R_{i_1,\dotsc,i_n}\cap C\neq\emptyset} |R'_{i_1,\dotsc,i_n}| \leq \sum_{R_{i_1,\dotsc,i_n}\cap C\neq\emptyset} \left(2KL/m\right)^n \leq \frac{(2KL)^n}{|R|}\sum_{R_{i_1,\dotsc,i_n}\cap C\neq\emptyset} \frac{|R|}{m^n} < \varepsilon\]
        \end{proof}
    \section{関数\texorpdfstring{$f,g$}{f,g}が\texorpdfstring{$\realNumbers$}{R}上で有界,連続で\texorpdfstring{$\integrate{-\infty}{\infty}{f(x)}{}{x},\integrate{-\infty}{\infty}{g(x)}{}{x}$}{f,gの-∞から+∞に渡る積分}が絶対収束するなら畳み込み\texorpdfstring{$f*g$}{f*g}は\texorpdfstring{$\realNumbers$}{R}上で連続で絶対可積分である。}
        \begin{proof}
            \quad\\
            (連続であること):
            \par
            $\varepsilon>0,x\in\realNumbers,h\in(-1,1)$を任意にとる。
            \begin{align*}
                \abs{(f*g)(x+h) - (f*g)(x)} &= \abs{\integrate{-\infty}{\infty}{\left[f(x-y+h) - f(x-y)\right]g(y)}{}{y}} \\
                &\leq \integrate{-\infty}{\infty}{\abs{\left[f(x-y+h) - f(x-y)\right]g(y)}}{}{y}
            \end{align*}
            $f$の有界性と$g$の絶対可積分性より，ある$a$が存在して
            \[ \integrate{-\infty}{a}{\abs{\left[f(x-y+h) - f(x-y)\right]g(y)}}{}{y},\integrate{a}{\infty}{\abs{\left[f(x-y+h) - f(x-y)\right]g(y)}}{}{y} < \varepsilon/3 \]
            となる。$f$は$\realNumbers$上で連続だから特に$[x-a-1,x+a+1]$上で一様連続である。
            よって$h$を十分小さくすれば
            \[ |f(x-y+h) - f(x-y)| < \frac{\varepsilon}{3L} \quad \left(L \coloneq \integrate{-\infty}{\infty}{|g(y)|}{}{y} \right) \]
            となり，
            \[ \integrate{-a}{a}{\abs{\left[f(x-y+h) - f(x-y)\right]g(y)}}{}{y} < \varepsilon/3 \]
            以上より
            \[ \abs{(f*g)(x+h) - (f*g)(x)} < \varepsilon \]
            \par
            \noindent
            (絶対可積分であること):
            \par
            $a>0$として$I(a) \coloneq \integrate{-a}{a}{|(f*g)(x)|}{}{x}$とすると$I(a)$は$a$に関して単調増加であり，次式が成り立つ。
            \[ I(a) = \integrate{-a}{a}{\abs{\integrate{-\infty}{\infty}{f(x-y)g(y)}{}{y}}}{}{x} \leq \integrate{-a}{a}{\integrate{-\infty}{\infty}{|f(x-y)||g(y)|}{}{y}}{}{x} = \integrate{-\infty}{\infty}{\integrate{-a}{a}{|f(x-y)||g(y)|}{}{x}}{}{y} \]
            右端の等号は$f$の有界性と$g$の絶対可積分性より$\integrate{-\infty}{\infty}{|f(x-y)||g(y)|}{}{y}$が$x$の値に関係なく収束することによる。
            これより
            \begin{align*}
                I(a) &\leq \integrate{-\infty}{\infty}{\abs{g(y)}\integrate{-a}{a}{\abs{f(x-y)}}{}{x}}{}{y} \leq \integrate{-\infty}{\infty}{\abs{g(y)}M_1}{}{y} \quad \left( M_1 \coloneq \integrate{-\infty}{\infty}{\abs{f(x-y)}}{}{x} \right) \\
                &\leq M_1M_2 \quad \left( M_2 \coloneq \integrate{-\infty}{\infty}{\abs{g(y)}}{}{y} \right)
            \end{align*}
            よって$I(a)$は単調増加かつ上に有界なので収束する。つまり$\integrate{-\infty}{\infty}{(f*g)(x)}{}{x}$は絶対収束する。
        \end{proof}
    \section{諸公式}
        \subsection{円座標表示された領域の面積\texorpdfstring{$\integrate{0}{2\pi}{r_x(\theta){r_y}'(\theta)}{}{\theta} = -\integrate{0}{2\pi}{r_y(\theta){r_x}'(\theta)}{}{\theta}$}{}}
            \label{円座標表示された領域の面積}
            領域 $S$ の境界が円座標表示 $x = r_x(\theta), y = r_y(\theta)$ で表されるとき、$S$ の面積は $\integrate{0}{2\pi}{r_x(\theta){r_y}'(\theta)}{}{\theta} = -\integrate{0}{2\pi}{r_y(\theta){r_x}'(\theta)}{}{\theta}$ である。
            証明は割愛する。
            左辺については，$\theta$ の増加に伴って，第 1 象限から第 4 象限にかけて順番に，横長の長方形を $y$ 軸方向に積み重ねてゆく様子を想像すれば納得できる。
            同様に，右辺についても長方形を積み重ねる向きを変えて理解できる。
        \subsection{\texorpdfstring{$n$}{n}次元の単体の体積\texorpdfstring{$\int_{x_1=0}^a \int_{x_2=0}^{a-x_1} \dotsi \int_{x_n=0}^{a-x_1- \dotsb - x_{n-2}} \mathrm{d}x_n \dotsi \mathrm{d}x_2\mathrm{d}x_1 = a^n/n!\;(a>0)$}{}}
            \begin{proof}
                \begin{align*}
                    &\phantom{=} \int_{x_1=0}^a \int_{x_2=0}^{a-x_1} \dotsi \int_{x_n=0}^{a-x_1- \dotsb - x_{n-2}} \mathrm{d}x_n \dotsi \mathrm{d}x_2\mathrm{d}x_1 \\
                    &= \int_{x_1=0}^a \int_{x_2=0}^{a-x_1} \dotsi \int_{x_{n-1}=0}^{a-x_1- \dotsb - x_{n-3}} (a-x_1- \dotsb - x_{n-1}) \mathrm{d}x_{n-1} \dotsi \mathrm{d}x_2\mathrm{d}x_1 \\
                    &= \frac{1}{2}\int_{x_1=0}^a \int_{x_2=0}^{a-x_1} \dotsi \int_{x_{n-2}=0}^{a-x_1-\dotsb-x_{n-4}} (a-x_1- \dotsb - x_{n-2})^2 \mathrm{d}x_{n-2} \dotsi \mathrm{d}x_2\mathrm{d}x_1 \\
                    &= \dotsb \\
                    &= \frac{1}{2\times 3\times\dotsb(n-1)}\int_{x_1=0}^a (a-x_1)^{n-1} \mathrm{d}x_1 \\
                    &= \frac{a^n}{n!}
                \end{align*}
            \end{proof}